\documentclass[14pt]{extarticle}
\usepackage{fontspec}
\usepackage{polyglossia}
\usepackage{amsmath}
\usepackage{amssymb}
\usepackage{xcolor}
\usepackage{fancyhdr}
\usepackage{graphicx}
\usepackage{listings}
\usepackage[most]{tcolorbox}
% \usepackage{setspace}
\usepackage{pifont}
\usepackage{enumitem}
\usepackage{titlesec}
\usepackage[bottom]{footmisc}
\usepackage{titling}
\usepackage{minted}
\usepackage{etoolbox}
\usepackage{array}
\usepackage{longtable} % page-breakable tables (must load before polyglossia/bidi RTL setup)
\usepackage{extsizes}
\usepackage{float}
\usepackage{booktabs}
\usepackage{hyperref}
\usepackage[onehalfspacing]{setspace}
\newif\ifwhiteblack
\ifdefined\setwhiteblack
  \setwhiteblack
\fi

\ifwhiteblack
  \usepackage[a4paper, landscape, top=1.25cm, bottom=1.25cm, left=1.25cm, right=1.25cm, headheight=1.5cm, headsep=0.5cm]{geometry}
\else
  \usepackage[a4paper, top=2.7cm, bottom=2.7cm, left=2.5cm, right=2.5cm, headheight=1.5cm, headsep=0.5cm]{geometry}
\fi
\usepackage{tikz}
\usepackage{datetime2}

\usetikzlibrary{automata, positioning, arrows.meta, arrows}
\providecommand{\LTcaptype}{table}

\newfontfamily\emoji{Segoe UI Emoji}

\pagestyle{fancy}

\setmainlanguage[numerals=western]{arabic}
\setotherlanguage{english}
\newfontfamily\arabicfont[Script=Arabic]{Calibri}
\newfontfamily\arabicfonttt[Script=Arabic]{Courier New}



\ifwhiteblack
\lstset{
  language=[Sharp]C,
  numbers=left,
  stepnumber=1,
  numbersep=8pt,
  frame=single,
  basicstyle=\ttfamily\small,
  keywordstyle=\bfseries,
  stringstyle=\ttfamily,
  commentstyle=\itshape
}
\else
\lstset{
  language=[Sharp]C,
  numbers=left,
  stepnumber=1,
  numbersep=8pt,
  frame=single,
  basicstyle=\ttfamily\small,
  keywordstyle=\color{blue},
  stringstyle=\color{red},
  commentstyle=\color{green!50!black}
}
\fi

\newif\ifdetailed
\ifdefined\setdetailed
  \setdetailed
\fi

\newif\ifwithsols
\ifdefined\setwithsols
  \setwithsols
\fi

\ifwhiteblack
  \AtBeginDocument{
    \LARGE
    \usemintedstyle{bw}
  }
\fi

% Default minted look: compact code that wraps long lines.
% Not applied in black-and-white landscape mode, which already uses \LARGE.
\ifwhiteblack\else
  \setminted{fontsize=\small, breaklines=true}
\fi

% unified tcolorboxes for articles
\ifwhiteblack
\tcbset{colback=white, colframe=black, fonttitle=\bfseries, boxrule=1pt}
\newtcolorbox{boxDef}[1][]{title={{\emoji📘} تعريف\ifx\\#1\\\else ~#1\fi :}}
\newtcolorbox{boxExercise}[1][]{title={{\emoji🧩} تمرين\ifx\\#1\\\else ~#1\fi :}}
\newtcolorbox{boxExample}[1][]{title={{\emoji📝} مثال\ifx\\#1\\\else ~#1\fi :}}
\newtcolorbox{boxNote}[1][]{title={{\emoji✨} ملاحظة\ifx\\#1\\\else ~#1\fi :}}
\newtcolorbox{boxAttention}[1][]{title={{\emoji🔔} تنبيه\ifx\\#1\\\else ~#1\fi :}}
\newtcolorbox{boxWarning}[1][]{title={{\emoji⚡} ملاحظة هامة\ifx\\#1\\\else ~#1\fi :}}
\newtcolorbox{boxSolution}[1][]{title={{\emoji✅} حل\ifx\\#1\\\else ~#1\fi :}}
\newtcolorbox{boxSymbol}[1][]{title={{\emoji🔣} رمز\ifx\\#1\\\else ~#1\fi :}}
\newtcolorbox{boxHint}[1][]{title={{\emoji💡} تلميح\ifx\\#1\\\else ~#1\fi :}}
\else
\tcbset{colback=white, colframe=black, fonttitle=\bfseries, boxrule=0.8pt}
\newtcolorbox{boxDef}[1][]{colback=blue!5!white,colframe=blue!75!black,
  title={{\emoji📘} تعريف\ifx\\#1\\\else ~#1\fi :}}
\newtcolorbox{boxExercise}[1][]{colback=cyan!5!white,colframe=cyan!70!black,
  title={{\emoji🧩} تمرين\ifx\\#1\\\else ~#1\fi :}}
\newtcolorbox{boxExample}[1][]{colback=yellow!5!white,colframe=orange!90!black,
  title={{\emoji📝} مثال\ifx\\#1\\\else ~#1\fi :}}
\newtcolorbox{boxNote}[1][]{colback=gray!10!white,colframe=black,
  title={{\emoji✨} ملاحظة\ifx\\#1\\\else ~#1\fi :}}
\newtcolorbox{boxAttention}[1][]{colback=magenta!10!white,colframe=magenta!80!black,
  title={{\emoji🔔} تنبيه\ifx\\#1\\\else ~#1\fi :}}
\newtcolorbox{boxWarning}[1][]{colback=red!5!white,colframe=red!75!black,
  title={{\emoji⚡} ملاحظة هامة\ifx\\#1\\\else ~#1\fi :}}
\newtcolorbox{boxSolution}[1][]{colback=green!5!white,colframe=green!60!black,
  title={{\emoji✅} حل\ifx\\#1\\\else ~#1\fi :}}
\newtcolorbox{boxSymbol}[1][]{colback=purple!5!white,colframe=purple!70!black,
  title={{\emoji🔣} رمز\ifx\\#1\\\else ~#1\fi :}}
\newtcolorbox{boxHint}[1][]{colback=teal!5!white,colframe=teal!60!black,
  title={{\emoji💡} تلميح\ifx\\#1\\\else ~#1\fi :}}
\fi


\ifwhiteblack
\tcbset{simplecode/.style={ colback=white, colframe=black, boxrule=1pt, arc=2pt, left=4pt,right=4pt,top=4pt,bottom=4pt}}
\else
\tcbset{simplecode/.style={ colback=gray!5, colframe=black!50, boxrule=0.4pt, arc=2pt, left=4pt,right=4pt,top=4pt,bottom=4pt}}
\fi
% Compact code boxes: tighter line spacing (the preamble uses onehalfspacing) and small vertical gaps.
% Must come AFTER the simplecode definition above.
\tcbset{simplecode/.append style={fontupper=\setstretch{0.95}, top=4pt, bottom=4pt, before skip=6pt, after skip=6pt}}
\newenvironment{boxCode}{\begin{tcolorbox}[simplecode]}{\end{tcolorbox}}

% Code cell used inside tables (keeps code left-to-right inside RTL tables).
% Use inside a fixed-width p{...} column: \begin{pycell}\begin{minted}...\end{minted}\end{pycell}
\newtcolorbox{pycell}{frame empty,colback=white,boxsep=0pt,left=2pt,right=2pt,top=2pt,bottom=2pt,width=\linewidth}

\newcolumntype{C}[1]{>{\centering\arraybackslash}p{#1}}

\makeatletter
\DTMnewdatestyle{dotted}{%
  \renewcommand{\DTMdisplaydate}[4]{\two@digits{##3}.\two@digits{##2}.##1}%
  \renewcommand{\DTMDisplaydate}{\DTMdisplaydate}%
}
\makeatother
\DTMsetdatestyle{dotted}

% redefine spaces after titles
\makeatletter
\renewcommand{\@maketitle}{%
  \begin{center}
    {\huge \bfseries \@title \par}%
    \vskip 0.2em % space between title and author
    {\large \@author \par}%
    % \vskip 0.2em % space between author and date
    % {\normalsize \@date \par}%
    {\small \hfill \begin{english}\DTMtoday\end{english} \par}%
  \end{center}
}
\makeatother

\fancyhf{} % clear default
\fancypagestyle{plain}{
  \fancyhf{}
  \fancyhead[L]{مدرسة التسامح الشاملة}
  % \fancyhead[L]{\includegraphics[height=1cm]{../../../images/logoTasamoh.png}}
  \fancyhead[R]{الأستاذ محمود اغبارية}
  \fancyfoot[C]{\thepage}
}

\fancyhead[L]{مدرسة التسامح الشاملة}
\fancyhead[R]{الأستاذ محمود اغبارية}
\fancyfoot[C]{\thepage}

\setcounter{tocdepth}{3} % only section subsection and subsubsection in TOC

\newcommand{\importpdfpage}[2]{%
    \thispagestyle{empty}
    \newgeometry{left=0cm,right=0cm,top=0cm,bottom=0cm}
    \includegraphics[page=#2,width=0.95\paperwidth,height=0.95\paperheight,keepaspectratio]{#1}
    \restoregeometry
}

\newcommand{\insertFullImg}[1]{%
    \noindent
    \makebox[\textwidth][c]{
        \includegraphics[width=0.9\paperwidth,keepaspectratio]{#1}%
    }%
}%


\ifwithsols
\title{حل اختبار شهري للصف العاشر}
\else
\title{اختبار شهري للصف العاشر}
\fi
\renewcommand{\DTMtoday}{29.10.2026}

\begin{document}
\maketitle
\thispagestyle{fancy}

\ifwithsols
\else
\begin{boxCode}
    \begin{itemize}
        \item يحتوي الاختبار على 4 أسئلة، \textbf{أجب عنها كلّها} $(4 \times 25 = 100)$.
        \item الوقت المخصص: ساعة ونصف.
        \item يسمح باستخدام كل مادة مساعدة، عدا الآلة التي يمكن برمجتها.
        \item اكتب بقلم حبر فقط، وبخط واضح.
        \item حلّ \textbf{السؤال الأوّل} على ورقة التسليم المرفقة في نهاية الامتحان، وحلّ باقي الأسئلة على ورقة الإجابة الخارجية (لا تنسَ كتابة اسمك الكامل).
        \item في كل الأسئلة افترض أن المدخلات صحيحة ولا داعي لفحصها إلا إذا ذُكر غير ذلك.
    \end{itemize}
\end{boxCode}

\vspace{1cm}
\begin{center}
  \includegraphics[width=0.4\textwidth]{../../../images/logoTasamoh.png}
\end{center}

\vspace{2cm}
\begin{flushleft}
    أرجو لكم التفوّق.

    الأستاذ محمود اغبارية.
\end{flushleft}
\clearpage
\fi

% ==============================================================
% السؤال الأول: فهم وتتبع كود معطى
% ==============================================================
\section{السؤال الأوّل}
أمامك مقطع البرنامج التالي الذي يستقبل عددين صحيحين $x$ و $y$:

\begin{boxCode}
\begin{english}
\begin{minted}{csharp}
int x = int.Parse(Console.ReadLine());
int y = int.Parse(Console.ReadLine());
int diff = Math.Abs(x - y);

if (x * y > 0)
{
    if (diff <= 5 && x > y)
    {
        Console.WriteLine("A");
    }
    else
    {
        Console.WriteLine("B");
    }
}
else
{
    if (x == 0 || y == 0)
    {
        Console.WriteLine("C");
    }
    else
    {
        Console.WriteLine("D");
    }
}
\end{minted}
\end{english}
\end{boxCode}

\begin{enumerate}[itemsep=2em]
    \filbreak
    \item \textbf{ما هو المُخرج؟} \\
    تتبّع تنفيذ المقطع البرمجي أعلاه لكل زوج من المدخلات في الجدول التالي، واكتب قيمة المتغير \textenglish{diff} والمُخرج الذي يطبعه البرنامج:
    \begin{center}
    \begin{english}
    \renewcommand{\arraystretch}{1.5}
    \begin{tabular}{|C{2.5cm}|C{2.5cm}|C{2.5cm}|C{3.5cm}|}
    \hline
    \textbf{$x$} & \textbf{$y$} & \textbf{diff} & \textarabic{\textbf{المُخرج}} \\
    \hline
    10 & 7 & & \\
    \hline
    -4 & -9 & & \\
    \hline
    5 & 8 & & \\
    \hline
    -5 & 3 & & \\
    \hline
    2 & 0 & & \\
    \hline
    \end{tabular}
    \end{english}
    \end{center}
    \filbreak

    \item \textbf{أعطِ مُدخلاً يطبع كلّ واحد من هذه المخرجات:} \\
    لكل مُخرج في الجدول التالي، اقترح زوجًا من الأعداد الصحيحة $x$ و $y$ (\textbf{مختلفًا} عن أزواج الفرع أ) يجعل البرنامج يطبع هذا المُخرج، مع تحقيق الشرط الإضافي المذكور:
    \begin{center}
    \renewcommand{\arraystretch}{1.5}
    \begin{tabular}{|C{2.5cm}|p{5.5cm}|C{2cm}|C{2cm}|}
    \hline
    \textbf{المُخرج المطلوب} & \textbf{شرط إضافي على المُدخلات} & \textbf{قيمة $x$} & \textbf{قيمة $y$} \\
    \hline
    $A$ & بحيث يكون العددان موجبَين & & \\
    \hline
    $B$ & بحيث يكون $diff \le 5$ & & \\
    \hline
    $B$ & بحيث يكون $x > y$ & & \\
    \hline
    $C$ & بدون شرط إضافي & & \\
    \hline
    $D$ & بحيث يكون $x$ موجبًا & & \\
    \hline
    \end{tabular}
    \end{center}
    \filbreak
\end{enumerate}

\ifwithsols
\begin{boxSolution}[الفرع أ والفرع ب]
\textbf{حل الفرع أ:}
\begin{center}
\begin{english}
\renewcommand{\arraystretch}{1.4}
\begin{tabular}{|C{2.5cm}|C{2.5cm}|C{2.5cm}|C{3.5cm}|}
\hline
\textbf{x} & \textbf{y} & \textbf{diff} & \textarabic{\textbf{المُخرج}} \\
\hline
10 & 7 & 3 & A \\
\hline
-4 & -9 & 5 & A \\
\hline
5 & 8 & 3 & B \\
\hline
-5 & 3 & 8 & D \\
\hline
2 & 0 & 2 & C \\
\hline
\end{tabular}
\end{english}
\end{center}

\textbf{حل الفرع ب (أمثلة لإجابات صحيحة):}
\begin{itemize}[nosep]
    \item لطباعة $A$ بعددين موجبين: يجب أن يكون $x > y$ والفرق $\le 5$، مثل: $x = 8,\ y = 5$.
    \item لطباعة $B$ مع $diff \le 5$: يجب أن يكون لهما نفس الإشارة ولكن $x \le y$، مثل: $x = 5,\ y = 8$.
    \item لطباعة $B$ مع $x > y$: يجب أن يكون لهما نفس الإشارة والفرق $> 5$، مثل: $x = 15,\ y = 4$.
    \item لطباعة $C$: يجب أن يكون أحدهما على الأقل صفرًا، مثل: $x = 0,\ y = 6$.
    \item لطباعة $D$ مع كون $x$ موجبًا: يجب أن يكون $y$ سالبًا، مثل: $x = 4,\ y = -2$.
\end{itemize}
\end{boxSolution}
\fi
\filbreak

% ==============================================================
% السؤال الثاني: تحليل عدد ثلاثي المنازل لخاناته
% ==============================================================
\section{السؤال الثاني}
نسمّي العدد الصحيح الموجب المكوّن من 3 منازل \textbf{«عددًا متوازنًا»} إذا كان \textbf{تربيع منزلة العشرات} (أي منزلة العشرات مضروبة بنفسها) مساويًا لـ \textbf{حاصل ضرب منزلة المئات في منزلة الآحاد}. \\
على سبيل المثال:
\begin{itemize}[nosep]
    \item العدد $248$ هو عدد متوازن لأنّ $4^2 = 16$ وكذلك $2 \times 8 = 16$.
    \item العدد $139$ هو عدد متوازن لأنّ $3^2 = 9$ وكذلك $1 \times 9 = 9$.
    \item العدد $352$ هو عدد غير متوازن لأنّ $5^2 = 25$ بينما $3 \times 2 = 6$.
    \item العدد $429$ هو عدد غير متوازن لأنّ $2^2 = 4$ بينما $4 \times 9 = 36$.
\end{itemize}

اكتب برنامجًا يستقبل من المستخدم عددًا صحيحًا موجبًا مكوّنًا من 3 منازل (بين 100 و 999):
\begin{itemize}[nosep]
    \item إذا كان العدد متوازنًا، يطبع البرنامج: \textenglish{"Balanced"}.
    \item وإلا، يطبع البرنامج: \textenglish{"Not Balanced"}.
\end{itemize}

\ifwithsols
\begin{boxSolution}
\begin{english}
\begin{minted}{csharp}
Console.WriteLine("Enter a 3-digit number:");
int num = int.Parse(Console.ReadLine());

int hundreds = num / 100;
int tens = (num / 10) % 10;
int units = num % 10;

int tensSquared = tens * tens;
int endsProduct = hundreds * units;

if (tensSquared == endsProduct)
{
    Console.WriteLine("Balanced");
}
else
{
    Console.WriteLine("Not Balanced");
}
\end{minted}
\end{english}
\end{boxSolution}
\fi
\filbreak

% ==============================================================
% السؤال الثالث: تعبير جبري ومجال التعريف
% ==============================================================
\section{السؤال الثالث}
معطى التعبير الجبري التالي الذي يعتمد على عددين عشريين $x$ و $y$:
$$ \text{result} = \frac{\sqrt{x - 4}}{y^2 - 9} $$

اكتب برنامجًا يستقبل من المستخدم العدد العشري $x$ ثم العدد العشري $y$. \\
قبل حساب قيمة التعبير، على البرنامج فحص مجال التعريف حسب الترتيب التالي:
\begin{itemize}[nosep]
    \item إذا كانت قيمة ما تحت الجذر التربيعي سالبة (أي عندما يكون $x < 4$)، لا يمكن حساب الجذر، فيطبع البرنامج: \\
    \textenglish{"Error: Negative root"}
    \item وإلا، إذا كانت قيمة المقام تساوي صفرًا (أي عندما يكون $y = 3$ أو $y = -3$)، لا يمكن القسمة على صفر، فيطبع البرنامج: \\
    \textenglish{"Error: Division by zero"}
    \item في جميع الحالات الأخرى، يحسب البرنامج قيمة التعبير $\text{result}$ ويطبعها.
\end{itemize}

\textbf{مثال:}
\begin{itemize}[nosep]
    \item إذا أدخل المستخدم $x = 2$ و $y = 3$، يطبع البرنامج: \\
    \textenglish{"Error: Negative root"}.
    \item إذا أدخل المستخدم $x = 20$ و $y = -3$، يطبع البرنامج: \\
    \textenglish{"Error: Division by zero"}.
    \item إذا أدخل المستخدم $x = 20$ و $y = 5$, يحسب البرنامج ويطبع: $0.25$، لأنّ: \\
    $\frac{\sqrt{20 - 4}}{5^2 - 9} = \frac{\sqrt{16}}{25 - 9} = \frac{4}{16} = 0.25$.
\end{itemize}

\ifwithsols
\begin{boxSolution}
\begin{english}
\begin{minted}{csharp}
Console.WriteLine("Enter x:");
double x = double.Parse(Console.ReadLine());
Console.WriteLine("Enter y:");
double y = double.Parse(Console.ReadLine());

if (x < 4)
{
    Console.WriteLine("Error: Negative root");
}
else if (y == 3 || y == -3)
{
    Console.WriteLine("Error: Division by zero");
}
else
{
    double numerator = Math.Sqrt(x - 4);
    double denominator = Math.Pow(y, 2) - 9;
    double result = numerator / denominator;
    Console.WriteLine(result);
}
\end{minted}
\end{english}
\end{boxSolution}
\fi
\filbreak

% ==============================================================
% السؤال الرابع: فاتورة المياه حسب حجم العائلة والاستهلاك
% ==============================================================
\section{السؤال الرابع}
تحسب شركة المياه فاتورة الاستهلاك الشهري للعائلات بناءً على \textbf{عدد أفراد العائلة} و\textbf{كمية المياه المستهلكة} بالمتر المكعّب. \\
اكتب برنامجًا يستقبل من المستخدم عدد أفراد العائلة (عدد صحيح موجب)، ثم كمية استهلاك المياه بالمتر المكعّب (عدد صحيح موجب)، ويحسب ويطبع مبلغ الفاتورة النهائي حسب القواعد التالية:
\begin{itemize}
    \item \textbf{إذا كانت العائلة صغيرة (أقل من $5$ أفراد):}
    \begin{itemize}[nosep]
        \item إذا كان الاستهلاك $20$ مترًا مكعّبًا أو أقل، تدفع العائلة $8$ شواكل عن كل متر مكعّب.
        \item إذا زاد الاستهلاك عن $20$ مترًا مكعّبًا، تدفع العائلة $8$ شواكل عن كل متر من أوّل $20$ مترًا، و $14$ شيكلاً عن كل متر مكعّب إضافي فوق الـ $20$.
    \end{itemize}
    \item \textbf{إذا كانت العائلة كبيرة ($5$ أفراد أو أكثر):}
    \begin{itemize}[nosep]
        \item إذا كان الاستهلاك $35$ مترًا مكعّبًا أو أقل، تدفع العائلة $7$ شواكل عن كل متر مكعّب.
        \item إذا زاد الاستهلاك عن $35$ مترًا مكعّبًا، تدفع العائلة $7$ شواكل عن كل متر من أوّل $35$ مترًا، و $12$ شيكلاً عن كل متر مكعّب إضافي فوق الـ $35$.
    \end{itemize}
\end{itemize}

\textbf{مثال:}
\begin{itemize}[nosep]
    \item إذا أدخل المستخدم عدد أفراد العائلة $3$ والاستهلاك $25$ مترًا مكعّبًا، يحسب البرنامج ويطبع: $230$، لأنّ:
    $20 \times 8 + (25 - 20) \times 14 = 160 + 70 = 230$.
    \item إذا أدخل المستخدم عدد أفراد العائلة $6$ والاستهلاك $30$ مترًا مكعّبًا، يحسب البرنامج ويطبع: $210$، لأنّ:
    $30 \times 7 = 210$.
\end{itemize}


\ifwithsols
\begin{boxSolution}
\begin{english}
\begin{minted}{csharp}
Console.WriteLine("Enter family members:");
int members = int.Parse(Console.ReadLine());
Console.WriteLine("Enter water consumption (cubes):");
int cubes = int.Parse(Console.ReadLine());

int bill;

if (members < 5)
{
    if (cubes <= 20)
    {
        bill = cubes * 8;
    }
    else
    {
        bill = 20 * 8 + (cubes - 20) * 14;
    }
}
else
{
    if (cubes <= 35)
    {
        bill = cubes * 7;
    }
    else
    {
        bill = 35 * 7 + (cubes - 35) * 12;
    }
}

Console.WriteLine(bill);
\end{minted}
\end{english}
\end{boxSolution}
\fi


% ==============================================================
% ورقة تسليم للطالب
% ==============================================================
\ifwithsols
\else
\clearpage
% \thispagestyle{empty}
% \mbox{}
% \clearpage

\begin{center}
{\Large \textbf{ورقة تسليم للطالب — السؤال الأوّل}}
\end{center}
\vspace{0.4cm}

\noindent\textbf{الاسم الكامل:}\hspace{0.4cm}\rule{8cm}{0.6pt}\hfill\textbf{الصف:} العاشر
\vspace{0.8cm}

\noindent\textbf{أ. ما هو مُخرج هذا المُدخل؟}
\begin{center}
\begin{english}
\renewcommand{\arraystretch}{2.1}
\begin{tabular}{|C{2.8cm}|C{2.8cm}|C{2.8cm}|C{4cm}|}
\hline
\textbf{$x$} & \textbf{$y$} & \textbf{diff} & \textarabic{\textbf{المُخرج}} \\
\hline
10 & 7 & & \\
\hline
-4 & -9 & & \\
\hline
5 & 8 & & \\
\hline
-5 & 3 & & \\
\hline
2 & 0 & & \\
\hline
\end{tabular}
\end{english}
\end{center}

\vspace{0.8cm}
\filbreak

\noindent\textbf{ب. أعطِ مُدخلاً يطبع كلّ واحد من هذه المخرجات (يجب ألّا يكون الزوج المختار مطابقًا لأي زوج في الفرع أ):}
\begin{center}
\renewcommand{\arraystretch}{2.1}
\begin{tabular}{|C{2.5cm}|p{5.8cm}|C{2cm}|C{2cm}|}
\hline
\textbf{المُخرج المطلوب} & \textbf{شرط إضافي على المُدخلات} & \textbf{قيمة $x$} & \textbf{قيمة $y$} \\
\hline
$A$ & بحيث يكون العددان موجبَين & & \\
\hline
$B$ & بحيث يكون $diff \le 5$ & & \\
\hline
$B$ & بحيث يكون $x > y$ & & \\
\hline
$C$ & بدون شرط إضافي & & \\
\hline
$D$ & بحيث يكون $x$ موجبًا & & \\
\hline
\end{tabular}
\end{center}
\fi

\end{document}
